\documentclass[11pt,a4paper]{article}
\usepackage[utf8]{inputenc}
\usepackage[T1]{fontenc}
\usepackage{amsmath,amssymb,amsthm}
\usepackage{booktabs}
\usepackage{algorithm2e}
\usepackage{hyperref}
\usepackage{geometry}
\usepackage{xcolor}
\usepackage{enumitem}
\geometry{margin=1in}
\hypersetup{colorlinks=true,linkcolor=blue!60!black,citecolor=green!40!black,urlcolor=blue!70!black}
\newtheorem{theorem}{Theorem}
\newtheorem{proposition}{Proposition}
\newtheorem{definition}{Definition}
\newtheorem{remark}{Remark}
\DeclareMathOperator{\tmark}{tanh}

\title{\textbf{GENESIS: Self-Improving Computation via Multiplicative Input Geometry}\\[6pt]
\large From Equation to Scaling Law in a Self-Modifying Dynamical System}
\author{Jun \\ SCI-Arc}
\date{February 2026}

\begin{document}
\maketitle

\begin{abstract}
We present a minimal self-improving computational system that uses one equation for both computation and self-modification. The core is a product-term dynamical map where multiplicative signal injection creates persistent attractor basin separation, enabling sequential memory without external loss functions. We prove by ablation that multiplicative and additive inputs serve structurally distinct roles that cannot be unified. We then show that online plasticity---where the system adapts its excitability parameters during computation based on its own signal sensitivity---consistently outperforms fixed-parameter systems across all tested dimensionalities ($D{=}12,16,20,24$), with the advantage ratio increasing from $+0.28$ at $D{=}12$ to $+0.91$ at $D{=}24$. The mechanism does not depend on scale to function but becomes more essential at scale. We document seven negative results that shaped the final architecture: offline restructuring accumulates noise, meta-plasticity self-throttles, dimensional expansion dilutes signal density, and selection cannot improve on birth alpha. The single positive path---online plasticity from the discrimination signal itself---is the only mechanism that survives all tests. We resolve a systematic measurement confound that caused apparent seed-dependent failures, and provide a GPU-accelerated implementation for testing at scales beyond pure Python.
\end{abstract}

\tableofcontents
\newpage

%===================================================================
\section{Introduction}
\label{sec:intro}
%===================================================================

The central question: \emph{what is the minimal mechanism by which a computational system can improve its own computation?}

This is not learning (which presupposes an external loss function). Not adaptation (which presupposes a designer who defines ``better''). It is the question of whether a dynamical system can, through its own operation, discover where its computation succeeds and use that discovery to restructure itself.

We approach this through increasingly constrained experiments, each building on the failures of the last. The constraints:
\begin{enumerate}[nosep]
    \item \textbf{Zero dependencies.} Pure arithmetic. No gradient libraries.
    \item \textbf{One equation.} The same map that computes must also provide the signal for self-modification.
    \item \textbf{No external evaluation.} The system must discover success from its own dynamics.
    \item \textbf{Honest measurement.} Every claim tested against random baselines with proper controls.
\end{enumerate}

The result is a product-term dynamical map with online plasticity that achieves consistent positive discrimination gaps across four dimensionalities, three birth seeds, and eight randomized signal worlds per condition---192 total tests.

\subsection{The Computational Task}

\begin{definition}[Sequential Discrimination]
Given $K$ signals $\{s_1, \ldots, s_K\} \subset \mathbb{R}^D$ and two orderings $\pi_1, \pi_2$ of $\{1,\ldots,K\}$, a system \emph{discriminates} the orderings if, after processing the signals in each order, the resulting attractor states are more similar within repeated runs of the same ordering than between different orderings.
\end{definition}

Let $c_{\pi,t}$ denote the centroid state after processing ordering $\pi$ on trial $t$. The \emph{discrimination gap}:
\begin{equation}
\label{eq:gap}
\mathrm{gap} = \underbrace{\left\langle \cos(c_{\pi,t_1}, c_{\pi,t_2}) \right\rangle_{\pi,\, t_1<t_2}}_{\text{within}} \;-\; \underbrace{\left\langle \cos(c_{\pi_1,t_1}, c_{\pi_2,t_2}) \right\rangle_{\pi_1\neq\pi_2,\, t}}_{\text{between}}
\end{equation}
Positive gap means the system retains order information.

%===================================================================
\section{The Core Equation}
\label{sec:equation}
%===================================================================

\subsection{Architecture}

The system: $N_C$ cells, each with state $x_i \in \mathbb{R}^D$ and per-cell excitability $\alpha_i \in \mathbb{R}^D$. The dynamics:

\begin{equation}
\label{eq:core}
\boxed{
\Phi_s(x_i)_k = \tmark\!\Big(\alpha_{i,k} \cdot x_{i,k} \;+\; \beta \cdot (x_{i,k{+}1} + \gamma s_{k{+}1})\,(x_{i,k{-}1} + \gamma s_{k{-}1})\Big)
}
\end{equation}

Indices are cyclic ($k{+}1 \equiv k{+}1 \bmod D$). Parameters: $\beta{=}0.5$ (product-term strength), $\gamma{=}0.9$ (signal injection). Signal $s \in \mathbb{R}^D$ or zero when absent.

State update with additive inter-cell coupling:
\begin{equation}
\label{eq:update}
x_i^{(t+1)} = (1{-}\delta)\,x_i^{(t)} + \delta\Big[\Phi_s(x_i^{(t)}) + \lambda_i \cdot \epsilon \sum_{j \neq i} w_{ij}\big(\Phi_0(x_j^{(t)}) - \Phi_0(x_i^{(t)})\big)\Big]
\end{equation}
where $\delta{=}0.35$, $\epsilon{=}0.15$, $w_{ij}$ are softmax attention weights, and:
\begin{equation}
\lambda_i = \exp\!\Big({-\|\Phi_0(x_i) - x_i\|^2}\,\big/\,{0.0225\,\|x_i\|^2}\Big)
\end{equation}
is a plasticity gate that closes during transient responses.

\subsection{Multiplicative vs.\ Additive: The Structural Distinction}

\begin{theorem}[Multiplicative-Additive Separation]
\label{thm:sep}
In Eq.~\eqref{eq:core}, expanding the product term yields:
\begin{align}
&(x_{k{+}1} + \gamma s_{k{+}1})(x_{k{-}1} + \gamma s_{k{-}1}) \notag\\
&\quad= x_{k{+}1}x_{k{-}1} + \gamma(s_{k{+}1}x_{k{-}1} + x_{k{+}1}s_{k{-}1}) + \gamma^2 s_{k{+}1}s_{k{-}1}
\end{align}
The term $\gamma^2 s_{k{+}1}s_{k{-}1}$ is state-independent and shifts the fixed-point landscape. Different signals create different basins. After signal removal, the system relaxes to a signal-determined basin: persistent memory.

Additive injection $+\gamma_{\mathrm{add}}s_k$ shifts the $\tmark$ argument uniformly. After removal, the shift disappears. No persistent memory.
\end{theorem}

\begin{table}[h]
\centering
\caption{Ablation: multiplicative vs.\ additive signal injection ($K{=}6$, 3 seeds).}
\label{tab:ablation}
\begin{tabular}{lccc}
\toprule
Architecture & Avg Gap & Pass Rate & Sequential Memory \\
\midrule
MULT (Eq.~\ref{eq:core}) & $+0.14$ & $7/9$ & \textbf{Yes} \\
ADD (additive signal) & $-0.02$ & $2/9$ & No \\
UNIFIED (single channel) & $-0.01$ & $3/9$ & No \\
\bottomrule
\end{tabular}
\end{table}

This was the first irrefutable result. The system achieved 11/12 criteria on a comprehensive test battery.

\subsection{Why Random Alpha Works}

For $\alpha_{i,k} \in [0.4, 1.8]$, the product term creates positive gaps with probability ${>}0.6$ across random draws. The cross-terms depend on signal geometry, not alpha. Alpha modulates response \emph{strength}, creating inter-cell diversity, but the memory mechanism operates through the product term regardless.

%===================================================================
\section{Capacity Analysis}
\label{sec:capacity}
%===================================================================

At $D{=}12$, $N_C{=}6$, testing $K \in \{3,\ldots,14\}$:

\textbf{Inverted capacity curve.} $K{=}16$ gap $(+0.14)$ exceeds $K{=}3$ gap $(+0.02)$. More signals can improve discrimination by creating richer attractor landscapes.

\textbf{Practical capacity.} $K/D \approx 0.50$ passes reliably. $K/D = 0.75$ passes 2/5 seeds. The bottleneck is alpha-signal alignment.

\textbf{Dimensional scaling (STILL, no plasticity):}

\begin{table}[h]
\centering
\caption{Architecture performance without plasticity.}
\label{tab:still}
\begin{tabular}{cccc}
\toprule
$D$ & Parameters & Avg Gap & Positive Rate \\
\midrule
12 & 72 & $+0.125$ & 75\% \\
16 & 96 & $+0.122$ & 71\% \\
20 & 120 & $+0.082$ & 67\% \\
24 & 144 & $+0.051$ & 58\% \\
\bottomrule
\end{tabular}
\end{table}

STILL performance degrades with $D$: random alpha covers less viable space in higher dimensions. This motivates the scaling test in Section~\ref{sec:scaling}.

%===================================================================
\section{Seven Failed Self-Construction Attempts}
\label{sec:failures}
%===================================================================

\subsection{GENESIS v1: Iterative Alpha Development (4/10 criteria)}

\textbf{Mechanism.} Uniform alpha $\to$ self-organize $\to$ expose to signals $\to$ measure product-term energy variance $\to$ crystallize alpha proportional to energy.

\textbf{Failure.} Energy variance captures state magnitude, not discrimination capacity. Development converged but did not outperform random.

\subsection{GENESIS v2: Self-Expanding Organism (8/10 criteria)}

\textbf{Mechanism.} Per-cell alpha. On failure: expand $D$ by adding dimensions. Grew $D{=}8 \to 32$.

\textbf{Failure.} RND-D8 dominated. Signals generated at $D{=}8$, padded with noise at $D{=}32$. Expansion diluted signal density.

\textbf{Lesson.} Product terms need signal density ($K/D$), not more dimensions.

\subsection{GENESIS v3: Self-Diagnosing Growth (3/10 criteria)}

\textbf{Mechanism.} Distinguished signal-sparse failure (restructure within $D$) from substrate-limited failure (expand $D$).

\textbf{Failure.} RESTR (blind restructure, fixed $D{=}12$) beat everything. Expansion was categorically wrong.

\textbf{Best result in project history:} RESTR-D12 achieved gaps $+0.31, +0.35, +0.25, +0.33$ at $K{=}8$--$11$.

\subsection{Seed v1: Offline Discrimination Restructuring (8/10 criteria)}

\textbf{Mechanism.} Run multiple orderings. Per-dimension: $|c_{\pi_1,k} - c_{\pi_2,k}|$ reveals discriminating dimensions. Per-cell: $\sum_k|x^{\pi_1}_{i,k} - x^{\pi_2}_{i,k}|/D$ reveals order-sensitive cells. Amplify alpha diversity on discriminating dims for discriminating cells.

\textbf{Result.} $+0.106$ over random. 3/4 seeds improved. But alpha moved only 0.04\% from birth.

\textbf{Lesson.} The discrimination signal is correct. But offline restructuring is a whisper.

\subsection{Compound: Meta-Plasticity (4/11 criteria)}

\textbf{Mechanism.} Level 1: online alpha plasticity. Level 2: $\eta$ adapts based on whether discrimination improved between epochs.

\textbf{Failure.} STILL won everything. Meta-adaptation correctly diagnosed plasticity was hurting and self-throttled $\eta$ from 0.0003 to 0.000052 (83\% reduction). Accumulated plasticity destroys.

\subsection{Selection: Evolutionary Approach (7/11 criteria)}

\textbf{Mechanism.} 4 offspring per generation, discrimination-informed perturbation, keep the best, 10 generations.

\textbf{Failure.} Beat random ($+0.041$) but not birth ($-0.168$). Evolved alpha specialized for curriculum endpoint, lost generality.

\subsection{Summary}

\begin{table}[h]
\centering
\caption{All approaches. Only online plasticity survives all tests.}
\label{tab:all}
\begin{tabular}{lccl}
\toprule
Approach & vs Birth & vs Random & Failure Mode \\
\midrule
Iterative dev (v1) & -- & -- & Wrong signal (energy) \\
Self-expansion (v2) & -- & -- & Dilutes density \\
Self-diagnosis (v3) & -- & -- & Expansion always wrong \\
Offline disc (seed v1) & $+0.037$ & $+0.106$ & Barely moves alpha \\
Meta-plasticity & $-0.115$ & $-0.115$ & Accumulation destroys \\
Selection & $-0.168$ & $+0.041$ & Overfits curriculum \\
\textbf{Online plasticity} & $\mathbf{+0.074}$ & $\mathbf{+0.074}$ & \textbf{Survives} \\
\bottomrule
\end{tabular}
\end{table}

%===================================================================
\section{Online Plasticity: The Surviving Mechanism}
\label{sec:online}
%===================================================================

\subsection{Algorithm}

After each step where signal is present:

\begin{algorithm}[H]
\SetAlgoLined
\caption{Online Plasticity (per step, cell $i$, dimension $k$)}
\label{alg:plast}
\KwIn{$\Phi_s(x_i)_k$, $\Phi_0(x_i)_k$, $\alpha_{i,k}$, rate $\eta$}
$r_{i,k} \leftarrow |\Phi_s(x_i)_k - \Phi_0(x_i)_k|$ \tcp*{signal sensitivity}
$z_{i,k} \leftarrow (r_{i,k} - \bar{r}) / \sigma_r$ \tcp*{z-score}
$\bar{\alpha}_k \leftarrow \frac{1}{N_C}\sum_j \alpha_{j,k}$; \quad $d_{i,k} \leftarrow \alpha_{i,k} - \bar{\alpha}_k$\;
\uIf{$|d_{i,k}| < 0.01$}{
    $\alpha_{i,k} \mathrel{+}= 0.3\eta \cdot \mathcal{N}(0,1)$ \tcp*{break symmetry}
}
\uElseIf{$z_{i,k} > 0$}{
    $\alpha_{i,k} \mathrel{+}= 0.5\eta \cdot \tmark(z_{i,k}) \cdot \mathrm{sign}(d_{i,k})$ \tcp*{amplify diversity}
}
\Else{
    $\alpha_{i,k} \mathrel{+}= 0.1\eta \cdot \mathcal{N}(0,1)$ \tcp*{gentle drift}
}
$\alpha_{i,k} \leftarrow \mathrm{clamp}(\alpha_{i,k},\; 0.3,\; 1.8)$\;
\end{algorithm}

$|\Phi_s - \Phi_0|$ measures how much the signal changed the dynamics at this cell and dimension. Dimensions with large effect and high inter-cell variance carry discrimination information. Alpha diversity on these dimensions is amplified.

\subsection{Why Online Works, Offline Fails}

Offline restructuring accumulates across epochs: 12 small pushes compound into overshoot. Online plasticity operates fresh each sequence---shifts during signal processing, then settle/final steps run without further plasticity. The shifts don't compound because each sequence starts from fresh initial state. Plasticity is \emph{transient}: helps current computation without permanently damaging the substrate.

The biological analogy: synaptic facilitation during perception (online), not long-term potentiation (offline accumulation).

\subsection{Three Uncertainties Resolved}

\begin{enumerate}[nosep]
\item \textbf{Restructuring vs.\ seed quality.} Restructured seed-42 outperformed birth seed-42 by $+0.037$ (7W/4L). Real.
\item \textbf{Discrimination signal vs.\ random perturbation.} Directed restructuring ($+0.182$) beat random perturbation ($+0.141$) by $+0.041$. Direction matters.
\item \textbf{Development metric alignment.} quick\_gap and measure\_gap agree on sign (4/4) but disagree on rank (2/4). Acceptable for guidance.
\end{enumerate}

%===================================================================
\section{The Seed-123 Confound}
\label{sec:confound}
%===================================================================

Across five artifacts, seed 123 appeared to catastrophically fail ($\delta = -0.186, -0.176, -0.155$). Diagnosis: all tests reused signal seeds derived from a fixed base, creating systematic alpha-signal alignment bias.

\textbf{Resolution.} With 8 randomized signal worlds per condition:

\begin{table}[h]
\centering
\caption{Performance with randomized signal worlds ($D{=}12$, $K{=}6$).}
\label{tab:confound}
\begin{tabular}{cccc}
\toprule
Seed & Avg Gap & Positive Rate & Range \\
\midrule
42 & $+0.117$ & 6/8 & $[-0.01, +0.31]$ \\
77 & $+0.131$ & 6/8 & $[-0.04, +0.43]$ \\
123 & $+0.121$ & 6/8 & $[-0.16, +0.41]$ \\
200 & $+0.126$ & 4/8 & $[-0.28, +0.48]$ \\
\bottomrule
\end{tabular}
\end{table}

Seed 123 is indistinguishable from others. No structural difference in alpha (mean 1.048 vs 1.052, std 0.407 vs 0.406, inter-cell cos 0.887 vs 0.882).

\textbf{Lesson.} Reusing signal generation seeds produces confounds. Randomization is essential.

%===================================================================
\section{The Scaling Result}
\label{sec:scaling}
%===================================================================

\subsection{Design}

ALIVE vs STILL at $D \in \{12, 16, 20, 24\}$. Per $D$: 3 birth seeds $\times$ 8 randomized signal worlds $=$ 24 tests. $K{=}6$, $\eta{=}0.0003$.

\subsection{Results}

\begin{table}[h]
\centering
\caption{ALIVE vs STILL across dimensions (24 tests per $D$).}
\label{tab:scaling}
\begin{tabular}{cccccc}
\toprule
$D$ & ALIVE & STILL & $\delta$ & Ratio & W/L \\
\midrule
12 & $+0.160$ & $+0.125$ & $+0.035$ & $+0.28$ & 15/7 \\
16 & $+0.148$ & $+0.122$ & $+0.026$ & $+0.21$ & 12/8 \\
20 & $+0.087$ & $+0.082$ & $+0.005$ & $+0.06$ & 14/10 \\
24 & $+0.097$ & $+0.051$ & $+0.046$ & $+0.91$ & 14/7 \\
\bottomrule
\end{tabular}
\end{table}

\textbf{Solid:} Delta positive at all $D$. Win rate exceeds loss rate everywhere. Online plasticity never hurts.

\textbf{Suggestive but not proven:} Delta trajectory $+0.035 \to +0.026 \to +0.005 \to +0.046$ is V-shaped, not monotonic. The ratio grows from $+0.28$ to $+0.91$ because STILL collapses at $D{=}24$, not because ALIVE improves. Four points with a valley is insufficient for a scaling law.

\textbf{Needed:} $D{=}32, 48, 64, 96, 128$ on GPU to resolve whether the V-shape is noise or structure.

\subsection{Per-Seed Trajectories}

\begin{table}[h]
\centering
\caption{Per-seed delta across $D$.}
\label{tab:perseed}
\begin{tabular}{cccccl}
\toprule
Seed & $D{=}12$ & $D{=}16$ & $D{=}20$ & $D{=}24$ & Trend \\
\midrule
42 & $+0.017$ & $-0.002$ & $+0.051$ & $+0.045$ & Grows \\
77 & $+0.143$ & $+0.048$ & $-0.007$ & $+0.029$ & Shrinks \\
200 & $-0.054$ & $+0.033$ & $-0.029$ & $+0.064$ & Grows \\
\bottomrule
\end{tabular}
\end{table}

2/3 seeds show growing delta. Seed 77 starts strong (its birth alpha is well-aligned at low $D$) and declines.

\subsection{Alpha Structure at Scale}

Alpha distance from birth: ${\sim}0.00006$--$0.00008$ at all $D$. Cell specialization: inter-cell cosine ${\sim}0.87$ at all $D$. The adaptation structure doesn't scale; its \emph{value} scales because the baseline (random alpha) degrades.

%===================================================================
\section{Theoretical Analysis}
\label{sec:theory}
%===================================================================

\subsection{Sequential Memory via Cross-Terms}

Signals $s_A, s_B$ in order $A \to B$: after processing $A$, state $x^{(A)}$ encodes $A$'s influence. Applying $B$:
$$\beta(x^{(A)}_{k{+}1} + \gamma s_{B,k{+}1})(x^{(A)}_{k{-}1} + \gamma s_{B,k{-}1})$$
creates $A \times B$ interaction terms distinct from $B \times A$ terms in reversed ordering. The $\tmark$ nonlinearity amplifies differences into distinct attractor basins.

\subsection{The Self-Referential Loop}

Eq.~\eqref{eq:core} simultaneously:
\begin{enumerate}[nosep]
\item \textbf{Computes:} cross-terms create basin separation.
\item \textbf{Reveals:} $|\Phi_s - \Phi_0|$ shows where signal has effect.
\item \textbf{Adapts:} online plasticity uses (2) to modify $\alpha$, improving (1).
\end{enumerate}

Same equation for computation and self-modification. No external loss. No separate evaluation. The computation \emph{is} the self-modification signal.

%===================================================================
\section{Discussion}
\label{sec:discussion}
%===================================================================

\subsection{What This Is}

A dynamical system with 72--144 parameters that discriminates sequential signal orderings, self-modifies during computation using its own sensitivity signal, consistently outperforms fixed parameters, works at small scale without dependence on large scale, and shows increasing advantage at larger scale (suggestive, not proven).

\subsection{What This Is Not}

Not intelligence (narrow task). Not consciousness (parameter update rule). Not a scaling law (4 points with V-shape). Not gradient descent (sign/magnitude heuristics, less efficient but self-contained). Not a singularity---at most a seed: the minimal mechanism that, given sufficient substrate, could compound.

\subsection{The Stability-Plasticity Resolution}

Small offline restructuring (${\sim}0.04\%$): noise. Large offline restructuring (12+ epochs): self-destruction. Online plasticity: transient, fresh each sequence, doesn't accumulate. The resolution is temporal: plasticity active during signal processing, inactive during settle/final.

\subsection{Open Questions}

Does the V-shape resolve at $D{=}32{+}$? Does $K/D \approx 0.50$ capacity hold at scale? Can the mechanism handle structured (non-Gaussian) inputs? Is there a principled way to derive $\eta$ from the dynamics?

%===================================================================
\section{Conclusion}
\label{sec:conclusion}
%===================================================================

One equation. The product term creates sequential memory through cross-terms. The sensitivity signal $|\Phi_s - \Phi_0|$ reveals where computation succeeds. Online plasticity amplifies alpha diversity on discriminating dimensions.

Works at $D{=}12$ ($\delta{=}+0.035$). Works at $D{=}24$ ($\delta{=}+0.046$). Wins majority at every tested $D$. Does not require scale. Seven alternatives failed. One survived.

The seed does not need the garden to germinate. But in richer soil, it grows taller.

\end{document}
